\section{실험 설계와 검증 질문}

문서군 전체를 다시 정리하면, 현재 논문에 필요한 실험은 방법 우열을 한 번에 선언하는 거대한 표가 아니다. 오히려 다음 세 질문에 대한 계층적 검증이다.

\begin{enumerate}
\item \textbf{구조 질문}: K-space는 정말 비등방적이며, 입력/코퍼스에 따라 의미 있는 facet 구조를 가지는가?
\item \textbf{회전 질문}: 정적 회전 선택 문제에서 PCA는 random/identity 및 다른 후보보다 더 나은가?
\item \textbf{시스템 질문}: 그렇게 얻은 회전 이득이 실제 low-bit PPL 우위로 이어지는가, 아니면 양자화기 설계에서 막히는가?
\end{enumerate}

\subsection{모델과 데이터}

구조 및 회전 질문의 주 실험장은 GPT-2 Medium이다 \citep{radford2019language}. GPT-2는 MHA 구조라 head별 공분산, 회전, 스펙트럼을 해석하기 쉽다. 시스템 질문에서는 Qwen 계열도 함께 본다 \citep{qwen2024qwen25}. 이는 현대 GQA 구조에서 baseline이 얼마나 강한지 확인하기 위함이다.

데이터는 두 층으로 나뉜다. calibration 단계에서는 technical, code, conversational 또는 WikiText-2 샘플로 공분산을 추정한다. 평가는 구조 실험에서는 KL, cosine similarity, rank correlation, MSE를 사용하고, 시스템 실험에서는 sliding-window perplexity를 사용한다.

\subsection{현재 논문이 실제로 사용하는 완료 실험}

현재 사실 기반으로 직접 재사용하는 실험은 네 가지다.

\begin{enumerate}
\item \textbf{Exp 2-1: facet separability}. 코퍼스 내 유사도와 코퍼스 간 유사도의 격차가 유의한지 본다.
\item \textbf{Exp 2-2: PCA vs LDA}. 공분산 정렬 회전이 지도 축보다 실제 양자화 surrogate에서 강한지 본다.
\item \textbf{Exp 2-3: query dynamicity}. importance schedule의 동적성이 축의 동적성까지 뜻하는지 분리해서 본다.
\item \textbf{Exp 3-1: continuous rotation optimization}. 단순 PCA가 실제로 강한 정적 기준선인지 확인한다.
\end{enumerate}

추가로 표준 PPL 비교는 Exp 4-2 로그를 사용하되, 이 논문에서는 그 수치를 \emph{중심 positive claim}이 아니라 \emph{어디서 병목이 생기는지 보여 주는 evidence}로 위치시킨다.

\subsection{가설과 판정 기준}

이 논문에서 직접 검증하는 가설은 다음 셋으로 제한한다.

\begin{description}[leftmargin=1.6em,style=nextline]
\item[가설 H1] K-space는 저차원 facet 구조를 가지며, 이는 단순한 노이즈가 아니다.
\item[가설 H2] 정적 회전 문제에서는 공분산 정렬 생성자(PCA)가 무작위 또는 의미론 기반 대안보다 강한 기준선이다.
\item[가설 H3] 현재 구현의 한계가 존재하더라도, PCA 축 자체의 유용성은 random/identity ablation과 surrogate 실험으로 분리해 확인할 수 있다.
\end{description}

반대로 \textbf{현재 논문이 아직 채택하지 않는 가설}은 ``FOKVQ가 모든 저비트 baseline을 이긴다''이다. 이 가설은 보류 상태로 두고, 계획 문서에서 별도 검증 대상으로 다룬다.
